% SOURCE docs/release/2026-10-29/source/docs/writing-repair/2026-10-29/src/c90569b7ab2e315b.tex CHAR 30
\chapter{Dataset choices and their causal roles}\label{pd15:datasets}}

A dataset is useful for the information it contributes to a decision, not simply because its rows concern finance. The comparisons proceed from funding panels and household records to servicing histories, experimental and generated benchmarks, then discovery, linkage and reasoning tasks. Each choice states a usable role and the information it cannot supply. The recommendations retain the survey's dated acquisition scope; links to subsequent studies distinguish fulfilled comparisons from remaining source and implementation work.

\hypertarget{funding-responses-and-bank-panels}{%
\section{Funding responses and bank
panels}\label{funding-responses-and-bank-panels}}

\textbf{1. FDIC Summary of Deposits.} The custodian describes an annual
survey of branch-office deposits as of 30 June. Institutions with
branches submit the survey; institutions with only a main office are
exempt from that reporting requirement. The natural observational unit
is a branch-year, nested within an institution and local market. This
makes the source attractive for deposit concentration, branch-network
histories and a proposed event-study workflow. Before estimation, the
chosen release must establish identifiers, ownership changes, deposit
units and coverage. A branch closure or merger is not automatically an
exogenous treatment: local demand, financial distress and strategic
selection can drive both the event and deposits. Annual June
observations cannot recover a daily customer offer response.
\textbf{Decision:} shortlist for a public panel prototype, with causal
event-study claims conditional on a separately justified design
\citep{pd15_fdic_sod}.

\textbf{2. FDIC Bank Data API financial and structural series.} The
inspected API specification exposes institution, financial,
historical-change and deposits routes. This gives an acquisition route
for an institution-time dataset and the structural history needed to
avoid treating a merger as organic growth. Financial fields and history
must be joined using a frozen field dictionary and defensible continuity
rules. First compare accounting identities, dates and missingness; then
examine a declared event or exposure. Bank and calendar dependence
belong in the uncertainty design. No bank-quarter dataset or causal
effect was downloaded or estimated in this survey. \textbf{Decision:}
pair with the deposits source after a bounded schema and vintage
inspection; do not pool institution, holding-company and branch units
without an explicit aggregation map \citep{pd15_fdic_api_spec}.

\textbf{3. Federal Reserve FR Y-9C.} The report collects consolidated
holding-company balance sheets, income statements and supporting
schedules quarterly. The official description says microdata are public
with exceptions. The unit is a consolidated holding company, which is
different from an insured bank or a branch. It can support funding
composition and institution-level stress examples when thresholds,
reporting changes and consolidations are accounted for. It is not a
customer-level dataset and does not supply the bank's internal pricing
decision. \textbf{Decision:} reserve as a second panel source; prefer
one coherent institution panel initially and add holding-company
aggregation only for a question that requires it \citep{pd15_y9c}.

\textbf{4. Federal Reserve H.8.} H.8 supplies estimated weekly aggregate
commercial-bank balance sheets and bank-group series, with downloadable
data routes. It is inexpensive context for system-wide deposits, credit
and common-calendar shocks. The aggregation is useful for a forecasting
or conditional-response example, but it removes customer and institution
assignment information. Group composition, revisions and seasonal
adjustments need attention. \textbf{Decision:} adopt as a contextual and
time-series workflow source; require a separately documented shock
design before attaching a structural causal interpretation
\citep{pd15_h8}.

\textbf{5. Archived macroeconomic extract.} The existing 203-quarter
dataset is already local and has an executed, bounded projection study.
Its advantages are a fixed historical snapshot, small resource
requirements and an existing independent numerical comparison. Its
limitations are age, aggregation and unresolved policy information.
\textbf{Decision:} reuse this baseline now; a newer series can be a
versioned extension, with release vintages and units preserved.
Replacing it merely because newer numbers exist would not supply an
identification argument \citep{public_macrodata_dynamic}.

\hypertarget{households-acquisition-and-credit-decisions}{%
\section{Households, acquisition and credit
decisions}\label{households-acquisition-and-credit-decisions}}

\textbf{6. SIPP-based 401(k) extract.} Eligibility, participation and
financial wealth yield a compact financial example with two distinct
estimands. The existing workflow can show adjustment-set declarations,
orthogonal scores, overlap diagnostics and why a participation effect
needs stronger assumptions than an eligibility contrast. Unknown survey
design and variable timing remain explicit. Employer selection and
saving preferences are plausible unmeasured confounders.
\textbf{Decision:} retain as the principal existing public
household-finance example; add sensitivity and target-definition tests
without relabeling observational assumptions as verified randomization
\citep{public401k}.

\textbf{7. Survey of Consumer Finances.} The 2022 public file is
described as 4,595 families represented by 22,975 rows across five
implicates. The Federal Reserve supplies separate replicate weights and
warns that ignoring both imputation and complex sampling yields
incorrect standard errors. This is valuable for wealth, debt,
balance-sheet heterogeneity and the product's treatment of uncertainty
across completed datasets. It is normally a triennial cross-section, not
a customer transaction panel. A credible exercise would reproduce a
published descriptive quantity with proper imputation and survey
variance before attempting an assumption-dependent causal example.
\textbf{Decision:} shortlist for observation construction and household
heterogeneity, selecting a documented pre-cutoff file; do not
concatenate implicates and count them as independent families
\citep{pd15_scf}.

\textbf{8. UCI Bank Marketing.} The data concern telephone campaigns at
a Portuguese bank and term-deposit subscription. The full additional
version has 20 inputs, 41,188 examples and a documented date ordering
from May 2008 to November 2010; the older full version has 45,211 rows.
The custodian states CC BY 4.0 terms. Contact duration is an especially
clear timing test: it is not available when deciding whether to place
the call. The files do not establish a random treatment policy, a
representative uncontacted group or an entity key sufficient for a
repeated-customer causal panel. ``The client subscribed'' is an outcome;
it is not itself an intervention. \textbf{Decision:} adopt immediately
for data profiling, business-question translation, leakage detection and
principled refusal. Reject a claimed incremental campaign effect from
this file alone \citep{pd15_uci_bank,pd15_profile}.

\textbf{9. UCI Default of Credit Card Clients.} This 30,000-row Taiwan
dataset has 23 predictors, including credit limits, payment status,
bills and payments over April--September 2005, followed by a default
label. The custodian states CC BY 4.0. The histories are useful for
checking feature clocks, censoring assumptions and the distinction
between a repayment and a bank action. A payment is partly a customer's
response, not an assigned servicing treatment. The file does not
disclose the complete sequence of contact, limit-change or relief
decisions and their logging probabilities. \textbf{Decision:} use as an
openly reusable credit-history and prediction baseline; use its
admissible covariates in a clearly labeled generated-action study if
desired. It cannot establish the causal effect of a collection policy
\citep{pd15_uci_credit}.

\textbf{10. Public HMDA application records.} The inspected CFPB
historical public schema includes lender identity, year, geography,
application action, loan amount, income and pricing fields. It supports
a lending-funnel example: application, denial, origination and
conditional loan terms are distinct objects. Conditioning a rate
analysis on origination changes the selected population. Public
disclosure transformations, missing fields and policy-specific coverage
require the selected year's dictionary. The current publication endpoint
returned HTTP 403; the accessible, cutoff-pinned historical schema is
evidence about that schema, not proof of the 2024 file's contents.
\textbf{Decision:} shortlist for application and selection analysis
after current-year access is resolved. Do not present it as a monthly
repayment panel or individual discrimination ground truth
\citep{pd15_hmda_schema}.

\hypertarget{loan-servicing-and-repeated-outcomes}{%
\section{Loan servicing and repeated
outcomes}\label{loan-servicing-and-repeated-outcomes}}

\textbf{11. Freddie Mac Single-Family Loan-Level Dataset.} The official
page describes origination and monthly performance data, including
prepayments and credit events, and directs users to register and sign
in. It also states that data are unaudited and subject to change. Free
noncommercial, academic/research and limited uses remain subject to
terms; commercial redistribution requires an agreement. Longitudinal
loan histories can demonstrate competing events, outcome horizons,
cohort splits and incomplete observations. Modification flags do not
supply the full decision state, all feasible offers or known action
probabilities. \textbf{Decision:} strong conditional candidate for
realistic servicing histories. Review the permitted use and a small
fixed cohort before committing resources; use an unrestricted synthetic
history route meanwhile \citep{pd15_freddie}.

\textbf{12. Fannie Mae Single-Family Loan Performance Data.} The
custodian describes selected fixed-rate loans with origination and
dynamic monthly performance, alongside a HARP mapping between original
and refinanced loans. Its stated population excludes several mortgage
categories, so it is not the entire mortgage market. Registration is
required; third-party redistribution and external commercial uses are
restricted without written consent. The HARP mapping is valuable for
discussing how an apparent loan termination may be a linked refinancing,
but it does not resolve endogenous refinancing selection.
\textbf{Decision:} alternative to Freddie Mac, not a mandatory second
acquisition. Select one source by access, variable coverage and cohort
quality; retain the other as a robustness extension only when it answers
a consequential question \citep{pd15_fannie}.

\textbf{13. NHEFS observational follow-up.} The inspected causallib
loader distinguishes a smoking-cessation/weight-change analysis and a
survival variant with follow-up and censoring. It explicitly records
optional exclusion of missing outcomes and uses cessation as the
exposure. This is an adjacent-domain teaching example for baseline
treatment, observation restrictions and survival semantics. It is not a
series of randomized bank servicing decisions. \textbf{Decision:} keep
as an optional inexpensive survival comparison after checking upstream
data terms and cohort construction. Use a generated sequential decision
process for time-varying treatment feedback rather than pretending this
baseline-exposure example supplies it \citep{pd15_nhefs_pinned}.

For a sequential policy, histories \(H_t\), actions \(A_t\), subsequent
states and censoring need a declared order. Identification by a
longitudinal g-formula requires consistency, positivity and sequential
exchangeability given the available history. Monthly rows satisfy none
of those assumptions merely by being longitudinal. An observed
modification may follow a private hardship assessment absent from the
release. Therefore the first servicing prototype should combine
realistic public outcome histories, where access permits, with generated
actions and outcomes whose causal law is fully declared. Its empirical
and generated components must remain separately labeled.

\hypertarget{experimental-panel-and-semi-synthetic-benchmarks}{%
\section{Experimental, panel and semi-synthetic
benchmarks}\label{experimental-panel-and-semi-synthetic-benchmarks}}

\textbf{14. NSW experimental job-training extract.} The
\texttt{causaldata} documentation describes \texttt{nsw\_mixtape} as a
445-row experimental extract with earnings and baseline covariates. A
small public earnings experiment is useful for effect-versus-prediction
semantics and a design-based baseline. Before claiming an identified
effect, recover the original assignment and sample-selection description
and bind the exact version. The original survey acquired a public teaching mirror without fitting a study. Chapter~\ref{gap29:public} subsequently uses the author's treated/control files and original selection argument for a fitted comparison, retaining their numerical difference from the mirror and the missing original assignment records.
\textbf{Decision:} preferred small experimental candidate, conditional
on those provenance and rights checks \citep{pd15_nsw_docs}.

\textbf{15. MatchIt LaLonde comparison sample.} The similarly named
\texttt{lalonde} file documented by MatchIt contains 614 observations:
185 treated observations from NSW and 429 PSID comparison observations.
It is not the all-randomized experimental comparison. This is a valuable
adverse example for an agent asked to ``load LaLonde and estimate the
RCT effect'': the response must disambiguate the dataset first.
\textbf{Decision:} retain as a selection-bias and dataset-identity
challenge, not the randomized baseline \citep{pd15_lalonde_pinned}.

\textbf{16. ACIC 2016.} The inspected author documentation specifies
4,802 empirical covariate rows and 58 covariates from the Collaborative
Perinatal Project. Its generator creates a treatment mechanism, response
surfaces and observations, and can return both potential outcomes,
conditional means and propensities. The package describes 77 settings
and 100 indexed seeds in its example; an evaluation must state the
chosen grid rather than quietly select easy settings. The code package
states GPL terms, but empirical input-data rights need their own review.
\textbf{Decision:} preferred standardized semi-synthetic estimation
benchmark. Freeze a modest scenario grid spanning overlap, nonlinearity
and heterogeneity, then compare a simple baseline and the proposed
method with identical inputs. Error against generated truth is
meaningful within that generator \citep{pd15_acic_pinned}.

\textbf{17. IHDP benchmark distributed with CEVAE.} The inspected
documentation explicitly says the benchmark contains generated treated
and control conditional means and factual/counterfactual outcomes. It
also warns that its legacy code is unmaintained. The name of an original
intervention study must not turn a generated benchmark into observed
counterfactual truth. \textbf{Decision:} secondary semi-synthetic
robustness option; prefer ACIC initially because adding a legacy
software stack contributes little to the first proposal. Read data in a
supported environment if this option is later used
\citep{pd15_cevae_pinned}.

\textbf{18. \texttt{did} county teen-employment panel.} The documented
\texttt{mpdta} example has 500 counties from 2003--2007, 2,500 rows and
a first treatment year associated with state minimum-wage changes. Its
\texttt{treat} indicator is ever-treated, not period-specific exposure.
This offers a compact staggered-adoption example and a concrete
treatment-clock test. Parallel trends, anticipation, state-level
assignment and limited assignment clusters need discussion; observed
adoption is not randomized. \textbf{Decision:} use as an adjacent-domain
panel adapter/reference example, with comparisons across appropriate
estimands. A future bank panel needs its own design \citep{pd15_mpdta}.

\textbf{19. Criteo uplift data.} The custodian describes advertising
incrementality tests in which a random subset is withheld from
targeting. Its page specifies 25 million users, 11 features, a treatment
indicator and visit/conversion labels. Feature projection and population
adjustments mean the original incrementality level cannot be recovered.
The displayed license is CC BY-NC-SA 4.0. \textbf{Decision:} useful
randomized large-data candidate, but deferred for this commercial
product proposal until permitted use is established. A free download is
not sufficient. A small generated randomized campaign with known
propensities can exercise the same interface without this dependency
\citep{pd15_criteo}.

\hypertarget{discovery-observation-construction-and-language-models}{%
\section{Discovery, observation construction and language
models}\label{discovery-observation-construction-and-language-models}}

\textbf{20. CausalBench perturbational gene-expression data.} The
documented benchmark offers observational, partially interventional and
fully interventional regimes on large biological datasets. It is
attractive because intervention labels allow evaluations that pure
observational graph recovery cannot supply. Nevertheless, perturbations
are biological operations, measured expression is an imperfect system
representation, and validation resources do not amount to a complete
known causal graph. Different auxiliary resources carry different
licenses, including a noncommercial component. \textbf{Decision:}
optional external-domain discovery challenge after a small declared
synthetic graph suite. Use held-out intervention prediction or
explicitly scoped reference metrics; do not call a learned biological
network banking graph truth \citep{pd15_causalbench_pinned}.

\textbf{21. FEBRL linkage datasets.} The inspected loader describes
synthetic identity records and returns known links reconstructed from
generator identities. FEBRL1 has 500 originals and 500 duplicates;
FEBRL4 has two tables of 5,000 records each with known correspondences.
These can test duplicate handling, blocking recall, false matches and
propagation of uncertain identity. Generator IDs that encode
correspondence must be hidden from matching features. A financial
outcome or transaction history added to these names is an additional
synthetic construction. \textbf{Decision:} preferred optional linkage
benchmark once data-specific terms are checked; existing
project-generated identities remain sufficient for an immediate
prototype. Neither establishes the error rate of a real bank audit
\citep{pd15_febrl_source}.

\textbf{22. PaySim.} The author description explains that mobile-money
transactions are generated by a simulator, initially from a service
schema and later using aggregate transaction information. Realistic
financial-looking records are still synthetic. This source is useful for
workload shape and event-stream concepts, but its fraud labels do not
identify the effect of a bank intervention. Code and linked dataset
terms need separate checks. \textbf{Decision:} optional engineering
fixture, deprioritized because the existing project simulator already
provides transparent decision and observation mechanisms
\citep{pd15_paysim_pinned}.

\textbf{23. CLadder.} The benchmark documentation describes 10,112
balanced generated causal questions, underlying model metadata and
variants with ordinary, counterintuitive and nonsense variable names.
Those variants help distinguish causal reasoning from memorized domain
associations. Public model metadata can support a logical reference; it
cannot be allowed into the evaluation prompt when it contains the
answer. \textbf{Decision:} shortlist for offline task design and later
bounded live evaluation after dataset terms and contamination controls
are fixed. Public benchmark performance is a diagnostic, not evidence
that the model owns the business assumptions
\citep{pd15_cladder_pinned}.

\textbf{24. Corr2Cause.} The documented generator builds graphs,
conditional-independence statements and verbalized causal-relation
questions; a three-class supplement distinguishes entailed,
underdetermined and contradicted claims within an equivalence class.
Variable renaming and paraphrases offer useful robustness contrasts.
\textbf{Decision:} pair selected tasks with CLadder and new controlled
banking briefs. Require the system to retain equivalent diagrams and
request missing identifying information instead of forcing a direction.
Source review is complete for these descriptions; no live-model scores
were produced \citep{pd15_corr2cause_pinned}.

The surveyed datasets cover concrete components of discovery, but none
establishes a universal ``business text to true causal graph'' route. A
useful graph prototype should return candidate structures, provenance
for asserted edges, explicit temporal restrictions and the consequences
of ambiguity for the requested intervention. On generated examples,
structural metrics can be scored against the generator. On empirical
observations, equivalence and uncertain measurement must survive into
the effect analysis.

The resulting portfolio favors one source per evidential role before adding scale. Existing household and macro examples support inexpensive scientific-interface comparisons; documented experiments support their own assignment contrasts; generated studies expose known failure mechanisms; and financial panels contribute measurement and continuity questions. Gated mortgage data, broader geographic coverage and uncertain reuse terms remain specific alternatives or omissions. None is needed merely to inflate the number of datasets, and none can replace a missing assignment or observation argument in another source.



% SOURCE docs/release/2026-10-29/source/docs/writing-repair/2026-10-29/src/0684ab9c1d754d1d.tex CHAR 0
\hypertarget{pd15:portfolio}{%

